\documentclass[11pt]{article}
\usepackage[margin=1.1in]{geometry}
\usepackage{amsmath,amssymb,amsthm}
\usepackage{booktabs}
\usepackage{hyperref}
\newtheorem{theorem}{Theorem}
\newtheorem{proposition}{Proposition}
\title{Simple symmetric Venn diagrams with 17 and 19 curves}
\author{Chris Dzoba}
\date{September 21, 2026 (draft)}
\begin{document}
\maketitle

\begin{abstract}
We exhibit simple, rotationally symmetric Venn diagrams with 17 curves and with 19 curves. Symmetric Venn diagrams
exist for every prime number of curves (Griggs, Killian and Savage, 2004), and simple ones were known only up to 13
curves (Mamakani and Ruskey, 2014). Four 17-curve and six 19-curve diagrams were found by a Metropolis walk on
rotation-invariant quadrangulations of the sphere in which regions may temporarily be duplicated, started from the
Griggs--Killian--Savage diagram with its multiple crossings resolved. Every diagram is given by a machine-checkable
certificate; one certificate of each size has been verified by a formal proof in Lean~4. All of the diagrams are
non-monotone, which is why the crossing-sequence searches that found the 11- and 13-curve diagrams could not have
found them.
\end{abstract}

\section{Statement}
An $n$-Venn diagram is a family of $n$ Jordan curves in the plane such that for every subset $S$ of the curves the
set of points inside exactly the curves of $S$ is nonempty and connected. It is \emph{simple} if no point lies on
three curves, and \emph{symmetric} if a rotation by $2\pi/n$ maps the family to itself. Gr\"unbaum asked in 1975
whether symmetric Venn diagrams exist for every prime $n$; Griggs, Killian and Savage answered yes in 2004, with
diagrams that are not simple. Simple symmetric diagrams were found for $n=11$ (Mamakani and Ruskey, 2012) and $n=13$
(Mamakani and Ruskey, 2014), and no larger case was known.

\begin{theorem}
Simple symmetric Venn diagrams with 17 curves exist, and so do simple symmetric Venn diagrams with 19 curves.
\end{theorem}

The proof is by exhibition. Four 17-curve diagrams and six 19-curve diagrams are given as certificates in the form
of Section~\ref{sec:encoding}, each accepted by two independently written checkers, and one of each size has been
formally verified in Lean~4 (Section~\ref{sec:lean}).

\section{Encoding and verification}\label{sec:encoding}
Let $D$ be a simple symmetric $n$-Venn diagram, drawn on the sphere with the unbounded region as a face. Its dual
map $M$ has one vertex per region, labelled by the region's membership pattern in $\{0,1\}^n$; one edge per arc,
joining labels that differ in exactly the bit of the arc's curve; and one face per crossing, a 4-cycle with labels
$u,\,u{\oplus}e_a,\,u{\oplus}e_a{\oplus}e_b,\,u{\oplus}e_b$ (a \emph{cube square}). Conversely:

\begin{proposition}
A map on the sphere whose vertices are labelled by distinct elements of $\{0,1\}^n$, all $2^n$ occurring, whose
faces are cube squares, whose edges of each coordinate $i$ form a single cycle through the faces containing them,
and whose vertex rotations are single cycles, is the dual of a simple $n$-Venn diagram; if the label rotation
$x\mapsto(x_1,\dots,x_{n-1},x_0)$ is an orientation-preserving automorphism of the map, the diagram is symmetric.
\end{proposition}

\begin{proof}[Proof sketch]
The map is a cellular embedding on the sphere; its geometric dual is a 4-regular plane graph whose vertices are the
faces (crossings). The edges of coordinate $i$, chained through opposite sides of each crossing, form one closed curve
$C_i$; regions of the arrangement are the vertices of $M$, and by the label condition each point of $\{0,1\}^n$ is
realized by exactly one region, which is a disk. Rotation of the labels acts as a rotation of the sphere fixing the
two vertices $0^n$ and $1^n$; taking $0^n$ as the outer face gives a plane diagram with $n$-fold rotational symmetry.
\end{proof}

A certificate is the list of the $2^n-2$ faces of $M$, each as four $n$-bit labels: $131{,}070$ faces for $n=17$
and $524{,}286$ for $n=19$. Verification recomputes, in seconds: all $2^n$ labels present exactly once, every edge in
exactly two faces ($2^{n+1}-4$ edges), Euler characteristic 2, every face transversal, a single rotation cycle at every
vertex, one cycle per curve with both sides connected, and invariance under label rotation. Two independently written
checkers, a half-edge implementation in C++ and a face-set implementation in Python, accept all ten certificates, as
does a third, written later and sharing no code with either, which audits the raw search state directly.

\begin{table}[h]
\centering\small
\begin{tabular}{llll}
\toprule
$n$ & certificate & SHA-256 prefix & pairwise distinct \\
\midrule
17 & \texttt{venn17-local-c3-s2} & \texttt{c178d7bd} & yes \\
17 & \texttt{venn17-gcp-s12} & \texttt{87f981b8} & yes \\
17 & \texttt{venn17-gcp-s14} & \texttt{eec8cd55} & yes \\
17 & \texttt{venn17-gcp-s16} & \texttt{431857aa} & yes \\
19 & \texttt{venn19-closure-s196002} & \texttt{fcb7ad34} & yes \\
19 & \texttt{venn19-closure-s196004} & \texttt{285de8ea} & yes \\
19 & \texttt{venn19-closure-s195001} & \texttt{a09ceae0} & yes \\
19 & \texttt{venn19-closure-s196001} & \texttt{56e7e1a1} & yes \\
19 & \texttt{venn19-closure-s196007} & \texttt{312e8c88} & yes \\
19 & \texttt{venn19-closure-s195002} & \texttt{726d1fb0} & yes \\
\bottomrule
\end{tabular}
\caption{The ten certificates. The six 19-curve face sets are pairwise distinct: the smallest symmetric difference
between two of them is $25{,}460$ faces (of $524{,}286$) and the largest $827{,}716$.}
\end{table}

\section{Formal verification}\label{sec:lean}
The first 17-curve certificate (\texttt{c178d7bd}) was formally verified in Lean~4 with Mathlib by Justin Grimes,
working from the definitions alone and without access to our code. His theorem states that the 17 curves
constructed from the certificate are Jordan curves in the Euclidean plane, carried to one another by a rigid
rotation through $2\pi/17$, with every one of the $2^{17}$ regions path-connected, no point on three curves, and every
crossing a transversal crossing of two curves. The finite parts of the proof use Lean's compiled evaluation
(\texttt{native\_decide}); the development builds from scratch in under an hour.

The first 19-curve certificate (\texttt{fcb7ad34}; JSON form \texttt{ed26b3ba}) was verified by a port of the same
development to $n=19$, carried out by Codex (OpenAI) from Grimes's sources: the dimension-dependent constants were
changed, the finite meridian and sector certificates that the proof consumes were regenerated for 19, and the
vendored topology library was left untouched. The port builds in twenty minutes and proves
\texttt{exists\_simple\_rotational\_venn\_19}: there exist $C:\mathrm{Fin}\,19\to\mathcal P(\mathbb R^2)$ and sides $S$
with \texttt{SimpleRotationalVenn}~$C$~$S$, the same predicate as at 17. Its axiom list is that of the 17-curve proof
under the namespace rename: \texttt{propext}, \texttt{Classical.choice}, \texttt{Quot.sound} and the twenty
\texttt{native\_decide} certificates. We re-ran \texttt{\#print axioms} on the built project to confirm this.

\section{Method}\label{sec:method}
We search the space of rotation-invariant cube-square quadrangulations of the sphere in which several vertices may
carry the same label. The energy of a state is $E = (\text{missing labels}) + \lambda\,(\text{duplicate regions})$;
$E=0$ with $\lambda>0$ is exactly a certificate. Moves are the local Reidemeister-type rewirings of the dual (lens
insertion and removal on a 3-cube, triangle flip, bigon insertion and removal), each applied simultaneously to all $n$
rotation images; every move preserves the sphere, the cube-square faces and the one-cycle-per-curve property, so the
walk never leaves the space of valid symmetric arrangements. Proposals are accepted by the Metropolis rule at
temperature $T$, with half of the proposals targeted at a missing label or a duplicated region.

Allowing duplicates is what makes the search work: with duplicates forbidden, greedy growth of the same
quadrangulations reaches $98.3\%$ of the labels and jams, because the last lens insertions require two missing labels
to coincide geometrically. With duplicates allowed, a missing label can be created next to a duplicate at zero cost and
defects diffuse until they annihilate.

\paragraph{17 curves.} The walk starts from the Griggs--Killian--Savage diagram for $n=17$ with every multiple
crossing resolved into simple crossings, a state with all labels present and $69{,}394$ duplicate regions. With
$\lambda$ ramped from $0.3$ to $1$ over three hours at $T=0.25$, three of fifteen independent runs reached $E=0$; a
fourth diagram was obtained by a further 40-minute cycle from a state with $68$ duplicates. Ten later instrumented
runs of the same recipe finished seven times within three hours, first reaching $E=2n$ after $1.6$--$2.8\times10^9$
proposals and finishing within $1.8\times10^9$ more. The same method finds 11- and 13-curve diagrams from the trivial
starting state in seconds and minutes.

\paragraph{19 curves: budget.} The resolved starting diagram for $n=19$ has all labels and $305{,}862$ duplicate
regions. Scaled by the number of regions, no 19-curve run in the first three days received the budget that the
17-curve runs needed: the longest had $6\times10^9$ proposals and was still descending when stopped, and every restart
from a low-energy state was shorter than the 17-curve endgame itself.

\paragraph{19 curves: the knot.} Every lineage that reached low energy stalled with the same duplicated label,
$\{0,3,4\}$, whose two instances had degrees 6--9 and shared a neighbour. A region of the dual can move by a triangle
flip only at degree 3, so the walk's targeted proposals, when they landed on such a duplicate, did nothing. The
duplicate is inherited: in the resolved starting diagram the label $\{0,3,4\}$ already appears twice, once as a region
of degree 18. (The resolved 17-curve diagram has the same feature for $\{0,1,4\}$, degrees 4 and 17, and the
17-curve walk happens to grind it down; the ten 17-curve pre-completion states all have their duplicates at degree 3
or 4.) One added proposal class removes the obstruction: when a targeted proposal lands on a duplicate of degree
above 3, it attacks that degree, by a triangle flip at a degree-3 neighbour or a lens insertion on an incident edge,
both ordinary moves under the ordinary acceptance rule. From a frozen $E=38$ state, 25 minutes at $\lambda=0.35$ with
this proposal reduced $\{0,3,4\}$ to multiplicity one; without it, 45 minutes at $\lambda=0.30$ left it untouched.
A three-hour ramp $\lambda: 0.5\to1$ then cooled the state to $E=38$ with two missing labels of rank 11 and no
duplicates.

\paragraph{19 curves: closure.} From that state, continuations differing only in the random seed were run with
$\lambda$ held at 1 and with $\lambda$ ramped from $0.6$, $0.8$ or $0.9$ to 1. Held at 1, five of eight
continuations reached $E=0$, after 4 seconds, 13 seconds, 190 seconds, 4.3 hours and 6.5 hours; of the sixteen
ramped continuations one finished, at the end of its six-hour ramp. The pre-completion states are saved exactly, by
journal undo of the winning move inside the engine, and show three closing mechanisms: a single triangle flip of the
last duplicated region, of degree 3, into the last missing label at Hamming distance exactly 3 (four diagrams); a
single lens insertion creating both missing labels at once after they had drifted to Hamming distance 1 (one
diagram); and a single lens removal of two adjacent duplicated regions in a state with no missing label (one diagram).

\section{Properties}
All ten diagrams are non-monotone: in the first 19-curve certificate, $3{,}590$ of the $27{,}594$ region orbits
($13\%$) lack a neighbouring region of rank one lower or one higher, and the 17-curve certificates behave the same
way. Every previously published simple symmetric Venn diagram with $n\ge 11$ is monotone, and the crossing-sequence
representation used to find them enumerates monotone diagrams by construction, so the diagrams reported here were
outside the search space of those methods. Minimum region degree is 3 (a degree-2 region is impossible in a simple
Venn diagram with $n\ge3$); in the first 19-curve certificate the degree histogram is $187{,}340$ regions of degree 3,
$199{,}120$ of degree 4, $98{,}686$ of degree 5, and a tail to degree 16, with the two poles of degree 19 and every
rank-1 and rank-18 region of degree 11.

\section{Data}
Certificates, pre-completion states, logs, the search code with its checksums, the independent checkers and both
Lean developments are at \url{https://github.com/dzoba/venn17} and archived on Zenodo (DOI to be inserted).

\section*{Contributions and use of AI systems}
This work was carried out by two AI systems working under the author's direction over five days, and this note
discloses their roles in full. Claude (Anthropic) formulated the dual-map encoding, designed and implemented the
search (the relaxed walk, its moves, energy and annealing schedules), diagnosed the failure at 19 from the run records,
designed and implemented the added proposal class and the capture of pre-completion states, ran the searches and the
verification, performed the analyses reported here and drafted this note. Codex (OpenAI) constructed the resolved
Griggs--Killian--Savage starting states, wrote independent validators for every certificate, identified the need for
the bigon moves, built exact local solvers used to diagnose the earlier plateau, and ported the Lean development to
19 curves. The author set the goals and strategy, decided what to pursue and what to abandon, supplied the computing
resources, and is responsible for the content.

\section*{Acknowledgments}
We thank Justin Grimes for the formal verification of the 17-curve certificate in Lean, carried out independently
and without knowledge of our construction, and on which the 19-curve verification is built; and Khalegh Mamakani for
his encouragement and for pointing out that the earlier searches excluded non-monotone diagrams by construction.
\end{document}
